%% tags: [mergesort]
%% source: 2023-fa-redemption_midterm_01

\begin{prob}
    Below is the \mintinline{python}{merge} function from the \mintinline{python}{mergesort} algorithm,
    with an added \mintinline{python}{print} statement:

    \begin{minted}[autogobble]{python}
        def merge(left, right, out):
            """Merge sorted arrays, store in out."""
            print(len(left))
            print(len(right))
            left.append(float('inf'))
            right.append(float('inf'))
            left_ix = 0
            right_ix = 0

            for ix in range(len(out)):
                if left[left_ix] < right[right_ix]:
                    out[ix] = left[left_ix]
                    left_ix += 1
                else:
                    out[ix] = right[right_ix]
                    right_ix += 1
    \end{minted}

    Suppose this version of \mintinline{python}{merge} is used in \mintinline{python}{mergesort}, and 
    \mintinline{python}{mergesort} is called on the following list:


    \mintinline{python}{[6, 2, 7, 1, 9, 3, 8, 5]}.

    If all of the numbers that are printed are added up, what will the sum be?

    \inlineresponsebox[.75in]{24}
\end{prob}
